\documentclass[9pt, a4paper]{article}
\usepackage[margin=0.5in, top=0.45in, bottom=0.45in]{geometry}
\usepackage{enumitem}
\usepackage{booktabs}
\usepackage{titlesec}
\usepackage[T1]{fontenc}
\usepackage{lmodern}
\usepackage{microtype}
\usepackage{xcolor}

\setlength{\parindent}{0pt}
\setlength{\parskip}{2pt}
\titlespacing*{\section}{0pt}{5pt}{2pt}
\titlespacing*{\subsection}{0pt}{4pt}{1pt}
\titleformat{\section}{\normalsize\bfseries\color{darkblue}}{}{}{}[\vspace{-2pt}\color{darkblue}\rule{\linewidth}{0.4pt}]
\titleformat{\subsection}{\small\bfseries\color{darkblue}}{}{}{}

\definecolor{darkblue}{RGB}{47, 84, 150}
\pagestyle{empty}

\begin{document}

\begin{center}
    {\Large\bfseries\color{darkblue} Rooftop Segmentation Using SAM}\\[2pt]
    {\small Zero-Shot Building Detection Across Three Indian Cities}\\[1pt]
    {\footnotesize\color{gray} AI/ML Internship Assignment -- Project 1 $\mid$ Resilience AI}
\end{center}

\section*{What This Project Does}

The goal was to automatically detect and outline rooftops from satellite images for three cities in India -- \textbf{Mumbai}, \textbf{Pune}, and \textbf{Vizag} -- without training any custom model. I used Meta's \textbf{Segment Anything Model (SAM)}, a powerful foundation model that can segment ``anything'' in an image. The key idea is: if SAM can segment general objects, can it find rooftops from a bird's-eye view? This project answers that question.

\section*{How It Works}

\begin{enumerate}[nosep, leftmargin=1.2em]
    \item \textbf{Download satellite images:} High-resolution tiles ($\sim$0.3--0.6\,m per pixel) are pulled from Esri World Imagery using the GeoJSON boundaries provided for each city.
    \item \textbf{Run SAM:} The largest SAM model (\texttt{vit\_h}) scans the image in a grid pattern, automatically detecting all distinct ``objects'' -- buildings, roads, trees, cars, everything. It processes the image in 1000$\times$1000 patches to avoid running out of GPU memory on Google Colab's free tier.
    \item \textbf{Filter out non-buildings:} SAM doesn't know what a ``building'' is -- it just finds shapes. So I applied three filters to keep only rooftops:
    \begin{itemize}[nosep, leftmargin=1em]
        \item \textit{Size filter:} Keep shapes between 15\,m$^2$ and 3000\,m$^2$ (removes tiny objects like cars and huge patches like fields).
        \item \textit{Shape filter (Rectangularity):} Buildings are roughly rectangular. If a shape's area is less than 60\% of its bounding box, it's probably a road or tree -- so it gets removed.
        \item \textit{Vegetation filter:} Uses the green channel from the RGB image to detect tree canopy and remove those false positives.
    \end{itemize}
    \item \textbf{Evaluate:} I compared my SAM detections against \textbf{Microsoft's Building Footprint dataset} (an independently generated ground truth) to compute IoU, Precision, Recall, and F1 scores.
\end{enumerate}

\subsection*{Bonus Experiment: Point-Prompted SAM}
I also tried a second approach: instead of letting SAM guess where buildings are, I fed it the \textit{center points} of known Microsoft buildings as hints. This gave much higher confidence per building ($\sim$0.7--0.88 vs $\sim$0.56), but obviously requires knowing where buildings are beforehand -- defeating the ``zero-shot'' purpose.

\section*{Results}

\vspace{2pt}
\begin{minipage}[t]{0.48\textwidth}
\centering
{\small\textbf{Pixel-Level Metrics (After Filtering)}}\\[3pt]
\footnotesize
\begin{tabular}{@{}lcccccc@{}}
\toprule
\textbf{City} & \textbf{SAM} & \textbf{GT} & \textbf{IoU} & \textbf{Prec.} & \textbf{Rec.} & \textbf{F1} \\
\midrule
Mumbai & 507 & 560 & 0.102 & 0.287 & 0.137 & 0.185 \\
Pune & 1,572 & 2,282 & 0.282 & 0.424 & 0.458 & 0.440 \\
Vizag & 486 & 630 & 0.327 & 0.465 & 0.524 & 0.493 \\
\bottomrule
\end{tabular}
\end{minipage}
\hfill
\begin{minipage}[t]{0.48\textwidth}
\centering
{\small\textbf{Instance-Level Matching (IoU $\geq$ 0.5)}}\\[3pt]
\footnotesize
\begin{tabular}{@{}lcccccc@{}}
\toprule
\textbf{City} & \textbf{TP} & \textbf{FP} & \textbf{FN} & \textbf{Prec.} & \textbf{Rec.} & \textbf{F1} \\
\midrule
Mumbai & 19 & 538 & 541 & 0.034 & 0.034 & 0.034 \\
Pune & 194 & 1,466 & 2,088 & 0.117 & 0.085 & 0.098 \\
Vizag & 134 & 353 & 496 & 0.275 & 0.213 & 0.240 \\
\bottomrule
\end{tabular}
\end{minipage}

\vspace{4pt}
\textbf{Vizag} performed best (F1 0.49) due to its well-spaced, structured layout. \textbf{Pune} was decent (F1 0.44) despite being very dense. \textbf{Mumbai} struggled badly (F1 0.19) because its tightly packed slums and high-rises have no visible gaps between buildings, so SAM merges entire city blocks into one giant blob.

\section*{Challenges I Faced}

\begin{itemize}[nosep, leftmargin=1.2em]
    \item \textbf{Building merging:} The biggest issue. When two buildings touch with no visible gap (common in Indian cities), SAM treats them as one object. This is why instance-level scores are much lower than pixel-level -- the total rooftop \textit{area} is roughly right, but individual buildings aren't separated.
    \item \textbf{Memory limits:} SAM's largest model is $\sim$2.5\,GB. On Colab's free GPU, I had to process images in small batches, aggressively free memory between cities, and delete the model immediately after use.
    \item \textbf{Tree canopy confusion:} SAM frequently detected clusters of trees as ``objects.'' The green-ratio vegetation filter fixed this -- removing 50 false positives in Mumbai and 88 in Pune.
\end{itemize}

\section*{What Could Be Done Next}

\begin{itemize}[nosep, leftmargin=1.2em]
    \item \textbf{Fine-tune SAM} on Indian rooftop datasets (e.g., SpaceNet) so it learns where one building ends and another begins.
    \item \textbf{Use multi-spectral imagery} (NIR bands) for perfect vegetation detection via NDVI instead of the approximate green-ratio method.
    \item \textbf{Hybrid approach:} Use a lightweight detector (UNet/YOLO) to find building centers first, then use SAM to trace the exact rooftop boundary for each one.
\end{itemize}

\end{document}
